\section{Results}
\subsection{Full-battery estimates and composition}
\label{sec:composition}
Full-battery margins exclude independence, while accuracy gives only a bound. Inflation reaches 1.244 on margins, where the raw interval excludes one, and 1.078 on accuracy, where it includes one (Table~\ref{tab:primary}). That upper endpoint maps under the stated measurement model and fixed battery to $\bar r_E<0.038$ and $K_{\mathrm{eff}}>7.47$.

\begin{table}[!htb]
\centering
\caption{Inflation estimates across 125 configurations with nominal 95\% recipe-cluster intervals. The auxiliary contrast compares the two auxiliary runs and removes the default-versus-auxiliary batch component, and the gain and competence adjustments are in Appendix~\ref{app:proxy}. The aggregate standard deviation $\sigma_{\mathrm{agg}}$ of the battery mean under the fitted covariance (Appendix~\ref{app:derivations}) uses margin units for margins and fractions for accuracy.}
\label{tab:primary}
\small\setlength{\tabcolsep}{4pt}\renewcommand{\arraystretch}{1.12}
\begin{tabular}{llrrrr}
\toprule
Score & Estimate & $\widehat\Lam$ & Interval & $K_{\mathrm{eff}}$ & $\sigma_{\mathrm{agg}}$\\
\midrule
Margin & Raw & 1.244 & \ci{1.143}{1.338} & 6.46 & 0.00572\\
Margin & Auxiliary contrast & 1.256 & \ci{1.090}{1.400} & 6.34 & 0.00562\\
Accuracy & Raw & 1.078 & \ci{0.993}{1.157} & 8.60 & 0.00514\\
Accuracy & Auxiliary contrast & 1.084 & \ci{0.939}{1.207} & 8.51 & 0.00516\\
\bottomrule
\end{tabular}
\end{table}

We observe that accuracy carries little seed signal at the benchmark level. The per-benchmark diagnostics that follow don't enter the headline ratio, which uses covariance products directly. Per-benchmark reliability averages 0.652 on margins and 0.324 on accuracy. Seven accuracy benchmarks and three margin benchmarks have reliability below 0.50, WinoGrande has a negative estimated variance on margins, and PIQA and WinoGrande have negative accuracy estimates. The accuracy correlation estimate also contains an entry of 1.296 and an eigenvalue of $-0.55$, and thus it isn't a valid correlation matrix without further regularisation.

We find that BoolQ dominates the battery, contributing 68\% of the margin covariance trace and 84\% of the accuracy trace. Removing it raises margin inflation to 1.786 with interval \ci{1.689}{1.877}, and accuracy inflation to 1.558 with interval \ci{1.412}{1.687}, while removing CommonsenseQA instead gives inflation estimates of 1.096 and 1.014. The removal estimates imply off-diagonal terms involving BoolQ summing to about $-0.15$ of the margin trace and $-0.06$ of the accuracy trace. Its removal therefore lowers the denominator of Equation~\ref{eq:lambda} more than the numerator, and lowers the aggregate standard deviation on both scales. A larger factor accompanies a less variable average. Lowering BoolQ's weight from a tenth to 0.0867 already stops the accuracy interval including one (Appendix~\ref{app:supplementary}). The factor compares each battery with independence at its own marginal variances.

Shared passages don't account for BoolQ's trace share, although that check can't reach other shared item effects. We recovered the 2,938 passages behind BoolQ's 3,270 questions from the release's request files, and assigning each passage whole to one half changes the two estimates by $-0.0002$ and $-0.0001$ on average across fifty re-splits. The BoolQ diagonal moves by 0.03\% and 0.15\% (Appendix~\ref{app:derivations}). A permutation that reshuffles each run's item deviations inside every group leaves all trait scores fixed and breaks the tie between particular items and runs. Across 2,000 such permutations 0.087 reach or exceed the observed margin estimate and 0.352 the observed accuracy estimate. Item-level sharing accounts for about 0.001 of the margin estimate and 0.002 of the accuracy estimate, but an effect carried by every item of a benchmark survives this permutation untouched.

Benchmark removal moves our estimates more than recipe or size removal. Recipe removal moves margin inflation between 1.207 and 1.256, and size removal between 1.178 and 1.289, where the deletion of 530M gives \ci{0.993}{1.338} and stops excluding one. The corresponding accuracy ranges run from 1.065 to 1.099 and from 1.047 to 1.115. Among equal-weight subsets of the same size, those containing BoolQ average 0.115 lower margin inflation at two benchmarks and 0.596 lower at nine than those without it. A sweep of 2,000 weightings from a flat Dirichlet leaves 0.817 of margin intervals excluding one. Figure~\ref{fig:composition} shows every single-benchmark removal with its interval on both scales, and Appendix~\ref{app:composition} separates benchmark count, identity, and weights. Our design treats this battery as fixed, and that spread isn't a distribution over benchmark collections.

\begin{figure}[t]
\centering
\includegraphics[width=0.95\linewidth]{figures/composition.pdf}
\caption{Inflation after removing each benchmark in turn, with nominal 95\% recipe-cluster intervals, from the trait-removal table of the extended supplement. The solid line marks one and the dashed line the full ten-benchmark estimate. On margins every removal leaves an interval that excludes one, and the estimate runs from 1.096 without CommonsenseQA to 1.786 without BoolQ.}
\label{fig:composition}
\end{figure}


\subsection{A registered test on unseen tasks}
\label{sec:heldout}
We find that the registered rule fails on held-out tasks. Held-out margin inflation on the 225 runs at 530M, 750M, and 1B is 1.217 with interval \ci{0.872}{1.498}, while the original ten on the same runs give 1.240 with interval \ci{1.054}{1.397} and pass it, a comparison we fixed before any held-out score existed. However, the two point estimates differ by 0.024, the held-out interval is 0.626 wide against 0.343 for the original, and our rule reads the lower limit. Interval width rather than the point gap decides the registered verdict.

Because the ten benchmarks above shaped our estimator, we registered this test on tasks outside them before any of their production scores existed. The protocol named SciQ, MedMCQA, and two AGIEval reasoning tasks, with a rule replacing any task whose 1B pilot accuracy fell below chance plus 0.02, a bar both AGIEval tasks missed by about a thousandth, inside the pilot's sampling error of about 0.010. The battery became SciQ, MedMCQA, and multiple-choice versions of DROP and CoQA, with 1,000 items in each of the four tasks and items sharing a passage or story kept in the same half. Our GPU budget couldn't cover all 375 released checkpoints, and our protocol's scope rule limited the test to those runs at the steps of the primary analysis. The test passes only if the lower 95\% limit exceeds one, and Table~\ref{tab:heldout} reports the test at its registered scope.

\begin{table}[!htb]
\centering
\caption{The registered test at 530M, 750M, and 1B, with nominal 95\% recipe-cluster intervals over 75 configurations. The auxiliary contrast compares the two auxiliary runs, the original rows read the released runs at the same sizes and were fixed before any held-out score existed, and the last row gives the median over 60 four-benchmark subsets of the original battery across all five sizes with their 5th and 95th percentiles, from the battery-size curve the protocol committed to.}
\label{tab:heldout}
\small\setlength{\tabcolsep}{4pt}\renewcommand{\arraystretch}{1.12}
\begin{tabular}{llrr}
\toprule
Battery & Score & $\widehat\Lam$ & Interval\\
\midrule
Held-out, four tasks & Margin & 1.217 & \ci{0.872}{1.498}\\
Held-out, four tasks & Accuracy & 1.220 & \ci{1.066}{1.359}\\
Held-out, auxiliary contrast & Margin & 1.181 & \ci{0.906}{1.410}\\
Original, ten benchmarks & Margin & 1.240 & \ci{1.054}{1.397}\\
Original, ten benchmarks & Accuracy & 1.097 & \ci{0.922}{1.246}\\
Original, auxiliary contrast & Margin & 1.273 & \ci{0.738}{1.629}\\
Original, four-benchmark subsets, all sizes (percentiles) & Margin & 1.176 & \ci{0.982}{1.476}\\
\bottomrule
\end{tabular}
\end{table}

Every run in scope finished scoring, the last 20 at 1B on three RTX 3090 cards after the weekly GPU quota ran out, and the scorer and settings match apart from torch 2.14.0 in place of 2.10.0, with our batch check passing by more than a factor of ten. In simulation at this scope with four traits, the rule passes in 0.926 of replicates at a true 1.244 under margin-like item noise and in 0.524 under accuracy-like noise. We added a noise level of 2.00 after seeing the two-size width, and its median width is 0.640, close to the observed held-out width. At that level the rule passes in 0.262 of replicates at a true 1.244 against a size of 0.025. On the same runs the accuracy interval is 0.293 wide. Three earlier results failed the rule as well, at 530M alone, at 530M and 750M, and at those two sizes with the 18 recipes first scored at 1B added, with estimates between 1.181 and 1.201 (Appendix~\ref{app:heldout}). The later power calculation can't rescue this registered test.

As an exploratory result, the held-out accuracy inflation interval excludes one. Inflation is 1.220 with interval \ci{1.066}{1.359}, above the 1.097 that the original ten give at the same scope. Removing CoQA lowers held-out margin inflation to 1.090, while removing any other task leaves it between 1.120 and 1.213, so the held-out estimate depends on battery composition as the original does (Appendix~\ref{app:heldout}).

